% ============================================================
%  Rasmus Nielsen, Philosophiske Grundproblemer (Philosophical Fundamental Problems)
%  Festskrift, University of Copenhagen 400th-anniversary, June 1879. 77 pp.
%  TRANSLATION (English)
% ============================================================
\documentclass[12pt]{article}
\usepackage[utf8]{inputenc}
\usepackage[T1]{fontenc}
\usepackage{libertinus}
\usepackage{libertinust1math}
\usepackage{textalpha}          % polytonic Greek typed directly (kept verbatim)
\usepackage[english]{babel}
\usepackage{enumitem}           % keep even if unused now — cheap insurance
\usepackage[protrusion=true,expansion=true]{microtype}
\usepackage[margin=1.4in]{geometry}
\usepackage{setspace}
\usepackage{fancyhdr}
\usepackage[colorlinks=true,linkcolor=black,urlcolor=black,
  pdftitle={Philosophical Fundamental Problems},
  pdfauthor={Rasmus Nielsen}]{hyperref}
\setlength{\headheight}{15pt}
\pagestyle{fancy}
\fancyhf{}
\fancyhead[LE,RO]{\thepage}
\fancyhead[RE]{\textit{Philosophical Fundamental Problems}}
\fancyhead[LO]{\textit{\rightmark}}
\setlength{\parindent}{1.5em}
\setlength{\parskip}{0pt}
\onehalfspacing

\title{\textbf{Philosophical Fundamental Problems}\\[0.5em]
  \large Festschrift for the University's Four-Hundredth Anniversary}
\author{R.\ Nielsen\\[4pt]\small Translated from the Danish}
\date{Copenhagen: Gyldendal, 1879}

\begin{document}
\maketitle


% ============================================================
%  INTRODUCTION (untitled)
% ============================================================
% Introduction (pp. 3--4)

% p. 3
\noindent
Philosophy does not, like other sciences, enjoy the advantage of
being able to build upon immediately firm presuppositions, upon
indisputable fundamental principles; for in the whole of philosophy
there is not a single proposition that stands so firm that it cannot,
when treated on its own, be assailed by criticism and dissolved in
dialectic. For that reason philosophy cannot very well parcel out its
proofs in the way that mathematics, for example, can. Although
mathematical proofs must be presented in strict order if full,
scientific insight is to be attained, the mathematician nevertheless
does not expose himself to any essential misunderstanding if, for one
purpose or another, he treats a particular problem in isolation. Those
who are competent know how the presuppositions can be proved, and those
whose prior knowledge does not suffice content themselves with the
result of the demonstration, trusting that the foundation is firm. But
upon such a trusting acquiescence the philosopher, it lies in the nature
of the case, dare not reckon. In philosophy no proof can stand alone;
the treatment of a separate part must borrow its light from the whole,
and the comprehensive system-proof must, so to speak, be transparent in
every single demonstration. But it is now well known from the history of
philosophy that the systems succeed one another, so that the preceding
ones always fall before those that follow. This is indeed to some extent
the case in the other sciences as well, but with the—admittedly very
significant—difference that certain discoveries, methods, and modes of
proof either remain unassailed or at least, under a further-reaching
elaboration, yield points of support, recognizable points of departure
for further progress. It is otherwise in philosophy. Here a new point of
view, a new thought, a new idea cannot step out into the world as a
physical or natural-historical discovery can. When a philosophical
system, with its principle and its peculiar method, falls prey to
critical dissolution and loses its standing, the whole is recast in such
a way that, in the new casting, one can hardly recognize what the abiding
element really is that the preceding system deposited. If one regards the
matter exclusively from this side, one would be tempted to believe that
philosophy was no true science. But the view is one-sided and therefore
misleading; the misleading element lies in the fact that one takes one's
models from other sciences and applies a standard that does not fit.
Philosophy must necessarily have its ideal peculiarity; what essentially
belongs to the invisible world of thoughts cannot possibly bear the same
marks as that which can be demonstrated from piecemeal-posited conditions
or exhibited in visible facts. Seen, on the contrary, from the standpoint
of the development of thought and of a consciousness growing in clarity,
it shows itself that philosophy, in steady historical continuation,
labors at the solution of certain ever-recurring problems; that every
thinker who, by placing one or more of these problems in a light
advantageous to inquiry, has made a contribution to the solution, remains
standing as an eminent figure in the history of philosophy; and that his
contribution continues to have effect, even though his school has long
since ceased to exist. For a closer understanding of this, and in order,
if possible, to contribute to a somewhat clearer conception of philosophy
as a science, attention is directed to the problems set forth below.


% ============================================================
%  I. THE PROBLEM OF KNOWLEDGE
% ============================================================
\section*{I.\ The Problem of Knowledge}
\label{sec:erkjendelse}

% p. 4 (section heading printed at the foot of p.4; text begins here)
\noindent
Among the philosophers who in more recent times have treated the
question of the conditions and limits of human knowledge, Imm.\ Kant
indisputably occupies the first place—if not by the results to which his
% p. 5
inquiries led him, then at any rate by the manner in which he has posed
the problem. Earlier thinkers—whether, like Locke, they laid the emphasis
on the senses, or, like Leibniz, on the understanding—simply presupposed
that we know from experience what experience is, that experience is
something that is understood of itself; but Kant's ``\emph{Kritik der
reinen Vernunft}'' shows that experience precisely is \emph{not} something
that is understood of itself, that in order to discover which cognitions
we can draw from the sources of experience, we must go back to the
inviolable fundamental conditions of cognition and begin with the question
of \emph{the possibility of experience}.

It is evident that the presuppositions which the knower himself must bring
with him in order to be able to make an experience cannot be borrowed from
experience, even though it is only with the help of experience that they
can come to consciousness. What, then, are these presuppositions? Without
sensation it is impossible to apprehend the sensible, and without
understanding impossible to determine what is sensed: experience is
conditioned by sensation and thought; but by what is sensation
conditioned, and how does it stand with the understanding? All sensation
rests upon an intuiting. To sense an object in space is just as impossible
without the intuition of space as it is impossible to apprehend a change
in time without the intuition of time. The pure intuitions, \emph{space
and time}, are therefore a priori necessary presuppositions for all
sensation and thereby for all experience. In order to procure a positive
content, thought must attach itself to sensation; for the understanding is
no intuiting faculty, and the old metaphysics' attempts to open up for
itself a supernatural world and, by pure understanding, to make
discoveries of supersensible things are fantastical. But that sensation
cannot bring it about that any experience whatever is made without the
assistance of the understanding is seen from this: that the one
sense-impression cannot be distinguished from the other without a
discriminating activity, and the one representation is not determined by
the other without a judging activity. The judging understanding would not
be able to exercise its activity if it did not itself contain
\textit{a priori} those root-concepts, categories, by which logical
judgments must necessarily be conditioned.
% p. 6
Which these categories are, and how many they are, can be seen from the
division of judgments in formal logic. Thus the hypothetical judgment
shows itself to be grounded in the categories of cause and effect, and the
disjunctive in the category of reciprocity. The a priori categories of the
understanding are therefore indispensable for all thought and thereby for
all experience.

The pure sense-intuitions and the a priori concepts must necessarily be
fundamentally different, indeed heterogeneous, since the senses are
heterogeneous with the understanding; the question then arises how it is
possible that forms so heterogeneous can be united in the cognizing
consciousness. We are referred to the imagination, which, on the one side
intellectual, on the other side sensible, by means of the intuition of
time, furnishes an a priori schema necessary for the possibility of
experience. It is not enough, for example, to have assured oneself of the
validity of the concept of cause as a concept alone; here a condition is
required for the concept's application to the appearances in the sensible
world, whereby the determinate order of time acquires decisive
significance.

By pointing out the subjective fundamental conditions holding for all
experience, Kant found himself able to overcome D.\ Hume's empirical
skepticism. Hume rejected the concept of cause as a kind of habit-concept
to which no experience corresponds. That things accompany and succeed one
another we can experience, but that the one is cause of the other we
cannot experience; what we by habit call cause is as unknown to us as what
we call force. Kant grants Hume this much, in that he demonstrates the
absurdity of wishing to derive the concept of cause from experience; but
inasmuch as he at the same time proves that no experience whatever is
possible without the concept of cause, he asserts victoriously and once
and for all this concept's a-priori-empirical validity so emphatically
that, according to him, experience and cause must be said to stand and
fall with one another. In particular cases of experience it can often be
doubtful what the cause is of an event that takes place. When A, B, and C
% p. 7
are accompanied by D, it may perhaps be A, perhaps B, perhaps C, perhaps A
and B, perhaps A and C, perhaps B and C, perhaps A, B, and C in
conjunction, in which the cause of D is to be sought; for the decision of
such doubtful cases one may employ induction, but induction precisely
presupposes that there must be a cause.\footnote{``Zwar scheint es, als
widerspreche dieses allen Bemerkungen, die man jederzeit über den Gang
unseres Verstandesgebrauchs gemacht hat, nach welchen wir nur allererst
durch die wahrgenommenen und verglichenen übereinstimmenden Folgen vieler
Begebenheiten auf vorhergehende Erscheinungen, eine Regel zu entdecken,
geleitet worden, der gemäss gewisse Begebenheiten auf gewisse
Erscheinungen jederzeit folgen, und dadurch zuerst veranlasst worden, uns
den Begriff von Ursache zu machen. Auf solchen Fuss würde dieser Begriff
bloss empirisch sein und die Regel, die er verschafft: dass Alles, was
geschieht, eine Ursache habe, würde eben so zufällig sein, als die
Erfahrung selbst; seine Allgemeinheit und Nothwendigkeit wären alsdenn nur
angedichtet und hätten keine wahre allgemeine Gültigkeit, weil sie nicht
\textit{a priori}, sondern nur auf Induction gegründet wären\ldots\
Freilich ist die logische Klarheit dieser Vorstellung einer, die Reihe der
Begebenheiten bestimmenden Regel, als eines Begriffs von Ursache, nur
alsdenn möglich, wenn wir davon in der Erfahrung Gebrauch gemacht haben;
aber eine Rücksicht auf dieselbe als Bedingung der synthetischen Einheit
der Erscheinungen in der Zeit, war doch der Grund der Erfahrung selbst und
ging also \textit{a priori} vor ihr vorher.'' Kritik der rein.\ Vernunft.
Leipzig 1838. S.~201.}

Here, be it noted, there is no question of the relation between the cause
and the thing itself, which falls outside our consciousness, but only of
the necessity of the concept of cause for our experience. An arbitrary
play of sensations and representations is only of subjective significance
and yields no experience. If a set of terms a, b, c, d, and so forth were
incessantly to alter their order and mutual position without any
conceivable ground whatever, we should indeed be compelled to apprehend
these changes in time, but the temporal succession would give us no
concept of cause and, as a matter of course, no experience either. If the
image of a house, for example, changed every moment in such a way that now
it turned its ridge downward, now upward, and, by the transposition of its
parts, took on the shape now of a tower, now of a whale, then we should
have to seek the ground of the changes in our own imagination and could in
no way refer such an image to any object of experience. For the sake of
experience there must be something that
% p. 8
compels us to place the one before and the other after; it is, in other
words, the necessarily determinate order of time that expresses the
relation between cause and effect. The necessity is a priori, for it is a
conceptual necessity; with this necessity the object is determined by
experience, and experience by the object.\footnote{``Der Begriff, der eine
Nothwendigkeit der synthetischen Einheit bei sich führt, kann nur ein
reiner Verstandesbegriff sein, der nicht in der Wahrnehmung liegt, und das
ist hier der Begriff des \emph{Verhältnisses der Ursache und Wirkung},
wovon die erstere die letztere in der Zeit, als die Folge, und nicht als
etwas, was blos in der Einbildung vorhergehen (oder gar überall nicht
wahrgenommen sein) könnte, bestimmt. Also ist nur dadurch, dass wir die
Folge der Erscheinungen, mithin alle Veränderung dem Gesetze der
Causalität unterwerfen, selbst Erfahrung d.\ i.\ empirisches Erkenntniss
von denselben möglich; mithin sind sie selbst als Gegenstände der
Erfahrung, nur nach eben dem Gesetze möglich.'' Anf.\ Skr.\ S.~196--97.}

From this one will understand that Kant, with inflexible rigor of
consequence, must make a divorce between the object of experience (the
phenomenon) and the thing in itself. There must be something that,
independently of us, gives our empty forms of cognition—the sensible as
well as the intellectual—a real content and determines our apprehension of
the object; but that the object, as we apprehend it, cannot possibly be
identical with the thing itself is a straightforward consequence of the
fact that, in passing through our subjective media of cognition, it must
necessarily take on a subjective tinge. Our apprehension of the thing is
an alteration of the thing: we see objects through colored spectacles.

The sharp separation between the thing in itself and its appearance finds
recognition in the exact sciences, particularly in the physiology of the
sense-organs. Helmholtz has\footnote{Handbuch der physiologischen Optik.
Leipzig 1867. S.~442--flgd.} emphatically stressed that our images and
representations are effects which the intuited and represented objects
produce upon our nervous system and thereby upon our consciousness. Here
there is a reciprocal action between an outside and an inside, and the
cause of the phenomenon is precisely this reciprocal action. To ask
whether the representation I have of a table—its form, its solidity, its
% p. 9
color, weight, and so forth—agrees with the real thing is the same as to
ask whether a tone is red, yellow, or blue. Things are determined by their
properties, and the properties by reciprocal action with our organs of
consciousness and with other things. Cinnabar is red, be it noted, for
those who can distinguish colors, but black for the red-blind. Although the
latter, according to the law of reciprocal action, have the same right in
their judgment as the former, we nevertheless say in general: cinnabar is
red; likewise we often ascribe to a thing a property as though it belonged
to the thing in itself, although this property holds only for the thing's
reciprocal action with other things, as when we say: lead is soluble, and
tacitly understand: in nitric acid, not in sulfuric acid. The result is
that our representations of things have only practical significance. Our
designations, our words and expressions, are only symbols; out of these
symbols a language can be formed, expressive of our observations, exact
investigations, and scientific cognitions, but which has nothing,
absolutely nothing, to do with the essence of things.

If this now is the final, decisive result of the inquiry, if the problem of
knowledge is herewith once and for all solved—in principle, that is—if all
further-reaching research will serve only for a further confirmation of the
Kantian theory: then not only must the works of Fichte, Schelling, and
Hegel be regarded as misguided, but every future philosophical attempt to
illuminate the essence and ground of existence must be dismissed as barren
and, for science, misleading. However the matter may stand, this much is
certain, that the Kantian doctrine of knowledge has an ambiguity that must
be drawn forth into the light. The fundamental question is: how can that
which fills our forms of cognition with \emph{real content be so absolutely
heterogeneous with these forms that the thing itself must remain entirely
unknown to us?}

Let us call the unknown X and confine ourselves to a definite object,
for example a table. That we apprehend the table in space and ascribe to
it a bounded extension, whereby it acquires its shape, is evident; but the
pure intuition of space gives
% p. 10
us no real table. The table is not a product of our own imagination; here
is a something that acts upon our senses, notably upon sight and touch—this
something is X. The table, the real phenomenon, thus comes about through
two heterogeneous contributions, a subjective one from ourselves and an
objective one from the thing in itself. We can symbolize it by an equation:
$X + Y = P_{x,y}$, where X denotes the objective, Y the subjective, and
$P_{x,y}$ a function of X and Y, the phenomenon. If we set $Y = 0$, then
the subjective vanishes, and $X = P_x$, that in the phenomenon which is due
to the objective alone, remains. That $P_x$ cannot itself be a phenomenon,
that subjectivity is required in order to make the thing into a phenomenon,
Kant has excellently shown; but how this happens, or even the possibility
that it can happen, he leaves standing in complete ambiguity. --- How,
then, does X actually stand in relation to $P_{x,y}$, the space filled by
the table? One can conceive the matter in several ways. X may be assumed,
without filling any space, to make an impression on the soul of such a
character that the determinate space-filling corresponds exactly to the
intensity of the impression. The real table is thus by no means a space
filled by X, for the X-impression brings it about only that subjectivity,
which itself fills the given space, must imagine that it is a real object
that fills it. This assumption must strike us as unreasonable, especially
when we consider that the table is not merely an object of sight, but also
an object of touch. Everywhere that we touch the table's surface, we sense
the reality. But the reality of every sensation derives, after all, from
X-impressions; there must then be just as many different X-impressions as
the table's surface has parts, and these impressions—whose difference here
is a difference of place—must be arranged and distributed in exactly the
same way as the parts of the surface, that is, arranged and distributed in
space. Yet neither does this assumption satisfy, for it would follow that
to the phenomenon of a square table there would correspond, if not indeed
a square X, then at least X-impressions that fit exactly the square, to the
phenomenon of a round table X-impressions that fit the round, and so forth.
If one separates the impressions
% p. 11
of X, which make the spatial determinations real, from X itself as that
which is outside space, then the X existing outside space has, after all,
nothing whatever to do with the real spatial phenomena; its significance
for a filling of the intuition of space is reduced to nil.\footnote{If one
assumes that X, in a certain negative way, reflects itself in its
X-impressions that refer to the intuition of space, just as the
X-impressions reflect themselves in the objects of experience, then one
gets an X-phenomenon behind every experiential phenomenon; but since the X
that stands behind the X-phenomenon must in turn reflect itself in an
X-phenomenon that removes itself still further from the form of space, one
gets X-phenomenon behind X-phenomenon ad infinitum before one reaches the
last X, separated from space, the pure ``Ding an sich''.}

Similar difficulties present themselves when we consider the phenomenal
changes that give the intuition of time real content. The
substance-phenomenon shows us what \emph{abides} in time, the
accident-phenomena what \emph{alternates} in time, and yet the substance
must be characterized by the accidents, and the accidents by the substance.
Now it is asked: is it one and the same X that determines the reality both
of the substance-phenomenon and of the accident-phenomena, or are there
here two kinds of X? If we assume the former, then we lay into X the very
same characteristic difference as into the phenomena. As the X of substance
it \emph{abides} in time, as the X of the accidents it \emph{alternates} in
time; this will then, in other words, mean that the thing in itself, the
unknown, is able to fill our intuition of time with real content only on
the inviolable condition that it, both from the side of substance and from
that of the accidents, enters into time, abides in time, and changes in
time. But this assumption stands in manifest conflict with the Kantian
presupposition. If we prefer the latter and posit one X for substance,
another for the accidents, then the multiplication of X will also be
carried further. It is then not enough to have a general X for substance in
general; we must have a particular X for each particular substance, a
sulfur-X for sulfur-substance, a
% p. 12
carbon-X for carbon-substance, an iron-X for iron-substance, and so on; as
regards the accidents, it will be doubtful how far we can make do with one
X for each group of accidents, or whether just as many X's are needed as
there are accidents in each group. The Kantian doubling—phenomenon and
thing in itself—comes, in this light, strongly to recall the Platonic: the
table and tableness, the horse and the horse-itself. However one turns and
twists it, the following question remains to be answered: Can X enter into
changes or can it not? If it cannot, then it keeps itself quite rightly
outside time, but then all phenomena of change become X-forsaken, unreal
semblances. If X can take part in the changes and give them the stamp of
experiential reality, then this reality-giving must also in one way or
another itself enter into time-determinations and take part in time. To
take up the X-impressions that determine the changes into the intuition of
time, and yet to let X be outside the intuition of time, is an absurdity,
for it is a contradiction.

The Kantian doctrine of knowledge impresses upon us that our concepts of
the understanding have validity only in the determination of the objects of
experience and find no application whatever to the thing in itself. But
when the thing in itself falls outside all conceptual determinations, how
then can it be determined? It is supposed not to be determinable by
sensation, since it is strictly held outside space and time; it must
therefore be determined by thought. Since now the understanding cannot
possibly procure for itself any positive content other than what it
receives from sensation, the consequence must surely be that the thing in
itself can be determined only negatively. One would now suppose that the
only determination that might be given of X would be this, that it excludes
all determinations. It would then be just as indeterminate that it is as
that it is not, just as indeterminate that it makes an impression on our
senses as that it makes no impression, and so on. But in this supposition
we are greatly disappointed. X becomes expressly determined by an act of
the understanding; it becomes determined by a judgment, not by a negative,
but by a positive judgment: \emph{the thing in itself is}. Being, a concept
of the understanding, is here,
% p. 13
contrary to the theory's own prescription, applied to X. Being denotes
reality, something positive; whence does the understanding get this
positive? Not from sensation, for sensation gives only what can be
apprehended in space and time. Nor from anything other than sensation, for
here there is nothing else that can give the understanding the positive. So
X must, then, be known indirectly by its effects, that is, by the
sense-impressions of which it is the cause. But this assumption too the
theory has excluded, for cause is a concept of the understanding just as
much as effect. By determining the thing in itself as an unknown cause
lying outside understanding and sensation—a cause that is to be thought as
objective and yet cannot be thought as objective, a cause that in its
effects, for the sake of reality, must be sensed and yet cannot be
sensed—Kant has, instead of showing the possibility of experience, on the
contrary proved the \emph{impossibility} of experience.

If there is no objective form at all that can belong to the thing in
itself, then this unknown must necessarily dissolve into a hazy
possibility, a ὕλη ἄμορφος, and lack all objective self-subsistence. It is
then so far from the case that the in-itself indeterminate possibility can
be said to determine the objective content of our subjective forms, that on
the contrary it borrows its apparent forms of reality from us and only
through this loan acquires the semblance of something objective. --- Here,
however, is a consideration that can be brought forward from the Kantian
side. Doubtless we must ascribe objective forms to the thing in itself,
insofar as we ascribe to it objective validity, but the point is that we do
not know these objective forms at all. To be sure, they cannot be forms
like ours; but who can prove that there are no other forms of existence
than those that underlie our sensation and thought? Not to speak of the
characteristic differences of our five senses—sight, hearing, smell, taste,
and touch—both among themselves and among different individuals, or of the
possibility that cognizing beings may have both more and sharper and finer
senses than we, there is surely no human being who can vouch that our a
priori intuitions and categories correspond even approximately to the
objective form-determinations expressed in
% p. 14
the thing itself. Is not the objectivity of space and time problematic?
Must we not, even as regards objects of experience, make a distinction
between the thing's concept and our imperfect, often faulty concept of the
thing? --- Here again we have the ambiguity.

It scarcely occurs to any reasonable person to deny the imperfection of our
intuition of space when it comes to a clear and distinct apprehension of
definite, many-sidedly combined and thereby involved relations in space. To
construct in space is not exactly an easy matter. Should anyone doubt this,
he need only make even a merely preliminary acquaintance with analytic
geometry; he will then soon learn how easily a curve or a mathematical
body, whose equation is known to him under certain conditions, can become
unrecognizable either by a displacement of the coordinates' point of
intersection or by the whole coordinate system's being altered. The more
combined the figure is, the more difficult are the special relations to see
through: here, then, manifold doubts can arise about the objectivity of
spatial relations in the particular. But it is one thing to concede the
imperfection of the representations and sensations conditioned by the
intuition of space, and another thing to deny objective validity to the a
priori form-determinations of the intuition of space. If a cognizing being
descended to us from higher regions and declared that it could distinguish
millions of parts in the space we took to be filled by a single part, we
should by no means contradict it; but if it moreover maintained that it
could distinctly see the circle as a square and the square as a circle,
should we then be convinced that its intuition of space was more objective
than ours? If the stranger gave proofs that he could not merely sense the
ether, but even count 733 trillion vibrations in one second, he would
astonish us in the very highest degree; but if he then gave as the reason
that he was able to apprehend the simultaneous as successive and the
successive as simultaneous, should we then dare to conclude from this that
his intuition of time came nearer to the objective than ours? Between space
and not-space, time and not-
% p. 15
time, no intermediate forms are possible. How many dimensions space may
appear to the sensing being to have is not here the question; the question
is of dimension and succession. By withdrawing itself from all
determinations of space and time, of dimension and succession, the whole
``objective, unknown world'' must precisely become known to us as a single
intensive point.\footnote{That the air-waves are heterogeneous with the
tones and the ether-waves with the colors proves only that a cause in the
objective can become unrecognizable by entering into reciprocal action with
the subjective; but what is incomprehensible herein proves nothing either
against the concept of cause or against the objective validity of time and
space. Of the influence of our organization there will be speech later.}

In confirmation that our concepts of reason (ideas) find just as little
application to anything objective in itself as our concepts of the
understanding, Kant sets up four so-called cosmological antinomies, namely
two mathematical and two dynamical. In this connection we confine ourselves
to the mathematical ones. First antinomy. \emph{Thesis}. The world has a
beginning in time, and is, as regards space, enclosed within limits.
\emph{Antithesis}. The world has no beginning, and no limits in space; it
is, in respect both of time and of space, infinite. Second antinomy.
\emph{Thesis}. Every composite substance consists of simple (uncomposite)
parts, and there exists nowhere anything but the simple and that which is
composed of it. \emph{Antithesis}. No composite thing in the world consists
of uncomposite parts, and there exists in general nothing simple in it.
--- Since the proof both for thesis and for antithesis in each of the
antinomies must be granted the same validity, this mutual conflict of the
transcendental ideas is supposed irrefutably to prove that we can just as
little reach the objective by means of our reason as by the use of our
understanding. Let us test the proofs. Does the thesis in the first
antinomy stand in insoluble conflict with the antithesis? We cannot see
that it does. If the world is from eternity, then there is, after all, no
contradiction whatever in the fact that up to a given point of time an
infinite series of successive world-states may have elapsed. To be sure, an
infinite series cannot by any successive synthesis
% p. 16
be completed in a finite time; it requires an infinite time in order to be
completed. But from whatever finite point of time one looks back toward the
infinitely distant beginning, one looks, after all, into an infinite time.
If there is nothing contradictory in assuming an infinite series of
represented terms, then there is nothing contradictory either in assuming
an infinite series of real terms, provided the reality, be it noted, is
presupposed to be from eternity. Why can the world not have begun in time?
Because, says the antithesis, a time must then have gone before, hence an
empty time; but in an empty time nothing can arise, since in no part of
such a time is there any condition for a distinguishing of existence from
non-being. But such an empty time is more imaginary than any imaginary
magnitude otherwise is. With an imaginary magnitude, mathematically
understood, one can nevertheless reckon; but an entirely empty time, a time
in which there is no difference between the past, the present, and the
future, is no time, and with no time one cannot reckon. As surely as the
real world really has a beginning, its beginning must necessarily fall in a
real time, and the real time begin with the real world. Since in this
respect it makes no difference whether we push the beginning's point of
time back half an eternity or a whole one, there is, as regards the world's
beginning, no conflict between thesis and antithesis. --- In considering
the limitation and non-limitation of space we come to the same result. The
real world is limited, for the outwardly real must at every moment have
definite limits and to that extent be a finite. But the world in its
finitude and limitation is also in real becoming, and can thus at every
moment be on the point of overstepping every finite limit. It is the real
world that makes space objective and sets objective limits in space; cosmic
space has no peculiar being outside the world. The antithetic is built on
the impossibility of completing the successive synthesis of the terms in an
infinite series, that is, on the impossibility of comprehending all parts
of an infinite totality. But, we ask, will then a limited
% p. 17
whole, which contains only 7 parts, be more real than a whole that contains
7 trillion parts, because the successive synthesis in the first case is
easy for us, but in the last case impossible? Surely no one will maintain
that. But if we may assume a totality with 7 trillion parts, a world-whole
with 7 trillion firmaments, to be real, notwithstanding that the successive
synthesis requires an enormously long time, how then can there be any
contradiction in presupposing the reality of an infinite totality, because
a successive synthesis (more correctly: coming-into-being) of its
individual parts must have required an infinitely long time?

For the \emph{Thesis} in the second antinomy it is adduced that nothing
substantial could be given if no indivisible parts existed; the
substantially composite must therefore consist of uncomposite parts. In
support of the \emph{Antithesis} it is adduced that no composition of
substances is possible except in space. The substance existing in space
must necessarily have substantial parts; but space has no substantial
parts, it is a \emph{totum}, no real \emph{compositum}. The contradiction
is now supposed to be: that composite substances cannot exist without
simple substances, and that simple substances cannot exist in space.
Instead of consisting of simple substances, the composite substance comes
to consist of composite substances ad infinitum, and that is absurd. ---
This contradiction, however, is more imaginary than real. If simple
substances and space are incompatible magnitudes, then not only experience
but also mathematics is impossible, for then intuition, and with it
sensation, is absolutely incompatible with thought. ``It is here,'' says
Kant, ``not enough, for the pure concept of the understanding of the
composite, to find the concept of the simple; one must find the intuition
of the simple that corresponds to the intuition of the composite, and this
is, according to the laws of sensibility, hence also in the case of objects
of sense, entirely impossible.'' --- This impossible, however, is not so
evident as it might seem at first glance. A mathematical body, a cube for
example, is an object of intuition; how does it stand here with the simple
and the
% p. 18
composite? The given cube can be intuited as a single existent, a continuum
bounded by 6 square faces; but it can also be regarded as a composite and
resolved into elements. If it is denoted by an infinite summation of
represented corporeality-elements $x^2dx$, then the third part of the
intuited cube is an integral, determined by finite limits, of these
elements. The face can be resolved into face-elements $xdx$, the line into
line-elements $dx$, and these again into elements of a higher order, and so
forth. In relation to the body the corporeality-element is the simple, just
as the face-element is the simple in relation to the face, and the
line-element the simple in relation to the line. Thought and intuition are,
as regards the mathematical body, agreed that the simple must be relative
for the sake of the composite, and the composite relative for the sake of
the simple, that therefore the fixing of an absolutely self-existent simple
is not an idea of reason, but a fixed idea. Should not the same hold for the
material body? It is supposed, according to Kant, to be an absurdity
``besides the mathematical point, which is simple, but no part, only a
limit in space, to think of physical points, which are indeed simple, but
yet have the advantage, as parts of space, of filling it by their
aggregation.'' But the alleged absurdity will not become evident to us. Let
us only put atoms and molecules instead of mathematical elements, and
aggregations instead of infinite summations: the fundamental relation
between the simple and the composite will nevertheless remain the same.
That the finite discretion is a parceling-out, a breach of continuity, must
be conceded. But even though the atoms are, in a real sense, ever so
indivisible—indeed even if, what for other reasons is an absurdity, they
are assumed to be absolute simple substances—the intuition of space can
nevertheless have not the least
% p. 19
to object against this. The atom's substantial simplicity by no means
stands in conflict with its mathematical divisibility. The extended-in-space
simple substance has, to be sure, just as many substance-elements as the
space it fills has space-elements, that is, it has infinitely many, for the
discretion goes on ad infinitum. But since the infinite discretion, as is
expressly seen from space, coincides with continuity, what can intuition
object against the simple substance—a real continuum—filling space, the
formal continuum? Since the parceling-out is presupposed, the finitude of
aggregation steps in place of the infinitude of summation; but since the
real distance between real parts precisely gives the form of space the
stamp of reality, the aggregation must surely be an eloquent proof of the
formal objectivity of space.

Kant is fully convinced of the insolubility of his antinomies, ``denn für
die Richtigkeit aller dieser Beweise verbürge ich mich\footnote{Prolegomena.
Werk.\ 3ter Band. S.~264.};'' but the insolubility presupposed, what is
thereby gained? ``That we must not think the phenomena of the sense-world
as things in themselves.'' That must be conceded. But there follows yet
something else from it, something that it is precarious to concede; for when
it is expressly added that we must not think the fundamental principles for
the connections of phenomena as valid for the things in themselves either,
then the understanding is lost, then it grows dark before the eye of reason.
If we could still stop at the preliminary result, that in all our cognizing
we must hold to our experiences and leave undecided how the things
themselves in the world unknown to us actually are constituted: that might
pass. But the antinomies drive us to an extreme that must either turn into
bottomless skepticism or awaken despair. We are to assume things in
themselves which, although they condition our experiences, nevertheless
exclude all conditions of experience, since neither space nor time, neither
unity nor plurality, neither ground nor consequence, neither cause nor laws,
neither whole nor parts, neither simplicity nor composition have the least
significance for these in-themselves unsensible and unthinkable things.
These are hard terms.

For the understanding of the conflict of the ideas of reason in the
antinomies, Kant adduces the following: ``Two mutually contradictory
propositions cannot both be false, unless the concept that lies
% p. 20
at the ground of them both is in itself contradictory; for example, the two
propositions: a square circle is round, and a square circle is not round,
are both false.'' But when the human consciousness has only a choice
between one of two: either to let the ideas of reason or the ``things in
themselves'' be ``square circles,'' then it is at least excusable if it once
more ventures to assert the right of reason by choosing the latter. This
choice is characteristic of J.\ G.\ Fichte's ``Wissenschaftslehre''.

We can briefly denote the deviation from the Kantian standpoint by
exchanging the symbol $X + Y = P_{x,y}$ for $XY = P$. That this must
necessarily give the inquiry another direction and lead to an entirely
different result is evident. When Y, the subjective, becomes nil, that is,
when the cognizing subjectivity, the Ego, vanishes, the whole becomes nil:
no indeterminable object, no self-existent X, remains. By establishing the
reciprocal determination of subjective and objective, Ego and non-Ego,
Fichte has introduced a great fundamental thought into the history of
philosophy, a fundamental principle of such incisive significance that it
can as little be excluded from the theory of knowledge as the law of
gravity from physics. Although Ego and non-Ego mutually exclude one
another, it is nevertheless self-evident that X, the non-Ego, can as little
be ascribed any self-subsistence valid outside the Ego as it can merge
together with the Ego itself. As real opposite to the Ego, the non-Ego must
therefore be an element in the Ego's essence, a barrier that is
indispensable for the Ego's self-development. Only by the thinking Ego's
becoming conscious of its opposite, its barrier, as an unconscious
presupposition that has proceeded from its own essence, does it become
conscious of itself, and overcome the barrier theoretically. The objects of
experience, the outward phenomena, present themselves immediately and with
a necessity over which the Ego has no command, and thereby take on a stamp
of outward reality; but that the Ego not only recognizes the forms and
categories in which the objects present themselves as its own forms, but
also recognizes the necessity with which
% p. 21
they present themselves as a necessity of the Ego's own nature-ground, is,
in the ``Doctrine of Science,'' the solution of the problem of knowledge.

The critical illusion in Kant consists in this, that the unknown thing
constantly appears as something self-subsistent in itself, objective,
positive—indeed as the only unconditionally positive—notwithstanding that
it shows itself from all sides that this positive can be determined only
negatively. The negative is positive, the positive is negative; this
contradiction, against which the critic defends himself with all his might,
crops up everywhere and vexes the Critique. This illusion Fichte has seen
through with clear consciousness, in that he has recognized the
significance of the negative and ascribed to X an essential value as
non-Ego.

But despite this unmistakable advance, the Fichtean solution is just as
unsatisfactory as the Kantian. What from the objective side is impossible
for Kant, namely to point out how the in-itself indeterminable X can fill
the forms of cognition with an unsurveyable manifold of real
determinations, the same is from the subjective side impossible for Fichte.
As non-Ego, X is far too uniform to yield any explanatory ground sufficient
for the diversity of the phenomena; a stone is a non-Ego, a plant a
non-Ego, an animal likewise: animals and plants have, in relation to the
Ego, the same significance as minerals and stones. The subjective
reflection of Ego and non-Ego consumes all objective relations between the
phenomena among themselves. Kant does at least get a world of experience,
even if on an ambiguous, self-dissolving foundation. Fichte, who removes
the ambiguity of the foundation, does not get even so much; he gets only a
subjectively necessary shadow-world of barriers and non-Ego, a world so
empty that the nature-content vanishes, and the philosophical interest in
the realities of experience is lost.

The Hegelian philosophy lets the problem of knowledge founder in the
absolute knowing instead of being solved. The opposition of the cognizing
subjectivity and the object of cognition, which is the bearer of the
problem, dissolves into the dialectical unity of thinking and being that
flows through the Logic
% p. 22
and issues in the Idea. In the Idea, the Absolute, subjectivity is
sublated, that is, it is transformed into a co-determination, a moment,
even if a moment of totality; and the same holds of objectivity, though
with a noticeable preponderance over the subjective. Instead of an
intuiting, thinking self-being that has the concept for its content and the
concept's reality for its object, we get in Hegel a concept that
comprehends itself, an Idea which, after having given itself subjectivity
in the form of reflection, again sublates this its subjectivity and
translates itself into the form of immediacy as objectivity in existing
objects. The concept, in absolute unity with its content \emph{the Idea},
cannot possibly maintain itself in two opposite forms at once; when the
Absolute is posited as subjective Idea in the form of subjectivity, that
is, of reflection, the objective is reduced to a moment and has vanished;
when the Absolute is posited as objective Idea in the form of objectivity,
that is, of immediacy, subjectivity has vanished: it must then comprehend
itself. But when the objective is to be able to comprehend itself just as
well as the subjective, then the self-comprehending concept is
incomprehensible. If the absolute Idea is to be the concluding unity of
subjective and objective Idea, then it must either assume a form which, in
being neither subjective nor objective, raises it above the extremes, or it
must be one of those two opposite forms in which it holds the extremes
conjoined. The former is, from Hegel's standpoint, inadmissible; for
Schelling's indifference of subjective and objective is only an abstract
presupposition, a non-ground, which stands far below the Ideas, and an
over-subjective or an over-objective has no warrant in ``the absolute
knowing.'' As regards the latter, we must then choose either the subjective
or the objective as the original form of the absolute unity; but the
Hegelian philosophy makes the choice impossible for us. If one compares the
Logic with the philosophy of nature, it appears as though the nature-Idea,
the Idea in the objective form of immediacy, were the derived, and the
Logic's Idea, the Idea in the subjectively reflected form, the original.
But by comparing the philosophy of nature with
% p. 23
the philosophy of spirit we are convinced that the matter stands the
reverse; for logic in logical reflection presupposes spirit. There can be
no logical reflection where there is no spirit; but nature is the
presupposition of spirit, just as immediacy is the presupposition of
reflection. Nature is spirit's ``An sich,'' and nature is the objective; it
is therefore the objective concept that—in an incomprehensible way!—
comprehends the subjective. That Hegel, by the immanent dialectic of the
ideal form-determinations, takes a giant step forward, that by going
further than Schelling he gets negativity developed by another measure and
in a quite different depth than either Kant or Fichte, must be conceded;
but the Absolute itself he does not get cleared up. Despite all the talk of
the self-thinking thinking, the absolutely self-knowing spirit, it is
impossible to learn what the knowing itself, the absolutely knowing,
is—whether it is absolutely subjective or absolutely objective, whether it
is a self or a self-less process or perhaps neither of the two. It is a
self-comprehending concept that is to clear up everything and bring the
last remnant of the unknown to vanish; but far from having eliminated X,
the incomprehensibly self-comprehending has transformed the Absolute itself
into an X, an impenetrable riddle that can fittingly be denoted by an
enigmatic symbol: $XY = X_{x,y}$.

If the problem of knowledge is to be solved, the following fundamental
conditions must inviolably be held fast:

\begin{enumerate}
\item[1)] \emph{There must be an insoluble opposition of subjective and
  objective, in which subjectivity itself, with its a priori forms, is the
  overreaching unity.}
\item[2)] \emph{Cognition must rest upon so thoroughgoing a reciprocal
  determination of the subjective and the objective that no object can
  either be or be thought to be without a corresponding objectivization.}
\item[3)] \emph{That imperfection of objectivization which has its ground in
  the independence of the existing objects from organically conditioned
  subjects must be capable of being remedied by subjective equivalents of
  ontological objectivization.}
\end{enumerate}

% p. 24
Against Fichte and Hegel, Kant is, as regards the first fundamental
thought, unmistakably right when he holds fast the opposition between the
cognizing subject and the object of cognition, just as he also—from the
side of finitude—is right in presupposing an unknown, a thing in itself,
that does not resolve into our cognition. But in the manner in which he
determines the thing in itself he is grossly wrong; for it is not by
excluding the a priori subjective form-determinations that the thing
withdraws itself from knowability\footnote{The more recent Critique has,
however, gone briskly forward, in that it has extended the influence of
``our organization'' to categories such as unity and plurality, substance
and accident, subject and predicate, and so on. We are, it is said,
naturally compelled to apply such categories on account of the peculiar
organization of our understanding; but it does not yet follow from this
that they are in themselves immutable. Enigmatic Critique! That the
understanding must not be conceived as a faculty by itself over against the
other faculties of the soul, that its activity must be dependent on organic
dispositions, can be seen from differing endowment; the understanding of
the individual may be disposed to one kind of intellectual activity without
being in the same degree disposed to another. Nor can it be denied that a
one-sided direction of the understanding may mislead one into fixing
fundamental concepts as exclusively a priori which yet are demonstrably not
so; one may think of the Pythagoreans' philosophical conception of the
decadic system, for example. But such missteps are, after all, to be
corrected by the understanding itself according to the understanding's
laws. Those critics, on the other hand, who give themselves the air of
being able to conceive—not merely an understanding infinitely higher than
the human, for that we can all do—no, an understanding that annuls the
difference between unity and manifold, cause and effect, subject and
predicate, must assuredly distinguish themselves by a quite peculiar
organization. ``Metaphysics'' they might now easily abolish; but how does
it go with the understanding? In order to conceive an understanding that is
absolutely heterogeneous with the human understanding, the human being
must, after all, take leave of his understanding.}, but by its apprehension
being modified according to the subjective peculiarity of the senses and of
the sensing beings. The unknown thing that is the object of sensation is a
thing with sensible properties just as well as the experiential phenomenon
that comes to our consciousness by way of sensation. In sense-systems of
different order, A, B, C, and so on, the same unknown X with its properties
must necessarily present itself in different phenomena. In the A-system the
properties $x_1, x_2, x_3, \ldots$ will present themselves
% p. 25
as $a_1, a_2, a_3\ldots$, in the B-system as $b_1, b_2, b_3\ldots$, and
thus in the first case give the experiential object $P_a$, in the last case
$P_b$. The cognizing beings have in the one system just as much right as
those in the other. But when the cognizing beings must all concede that the
thing itself falls outside their cognition, on what then does the truth of
cognition rest? The answer to this question leads us to the second
fundamental condition.

If the object in itself, with its properties, its concrete determinations,
could really stand independently outside all objectivization, and if the
objective thus standing were the really true, then the cognition of truth
would be impossible. But by working through the relation between Ego and
non-Ego, Fichte has strikingly proved that the objective, far from being
independent, is dependent on and determined by subjectivity. The
exaggeration with which he asserts the pure Ego's original self-limitation,
and the energy with which it itself posits its objects, furnishes a most
necessary counterweight against the vulgar assumption that objectivization
does not really concern the objects, but that it becomes subjectivity's own
affair whether it objectivizes truly or falsely. When we now concede to
Fichte that every real object must be a necessary product of an
objectivizing subjectivity, but cannot possibly concede to him that the
objects really are in themselves as we objectivize them to ourselves, or
that our human objectivization should be able to give the things existence,
then we resort to the third fundamental condition, in order to discover how
the imperfection of the finite objectivization is to be remedied.

The universally valid ground of cognition coincides with the ontological
ground of objectivization. When someone, for example, propounds the
hypothesis of the rotating primordial vapor out of which our planetary
system is supposed to have developed, he certainly does not assume that
that vapor was an object of experience for human sensation, any more than
his opinion is that what was then ``a thing in itself'' has since, through
our experiences, been transformed into vapor. Nor does he assume either
that the primordial vapor
% p. 26
existed in our subjective intuition of space and rotated in our subjective
intuition of time, nor that it is we who with a priori necessity have
prescribed to nature the laws according to which the whole development has
taken place. No, he rather assumes that all of it—the space, the time, the
vapor, the motion—was objective without any objectivization whatever. Since
it is now proved that the last assumption is just as meaningless as the
first, what follows from this? That here in silence there must have been
slipped in a subjective equivalent of ontological objectivization!
Hypothetically this is expressed thus: if we had been present and been able
to survey the whole, then the matter must necessarily have appeared as the
description runs. Thus: every time our limited subjectivity goes beyond its
own bounds of experience, it becomes, for the sake of the objective,
compelled to presuppose an intuiting that is not our intuiting, a seeing
that is not our seeing, a thinking that is not our thinking, an
objectivization that is not our objectivization; it becomes, in short,
compelled to presuppose an ontologically objectivizing subjectivity.

From the ontological fundamental law of objectivization, symbolized by
$XY = O$, where O denotes the object, it follows: that every objective
moment must necessarily presuppose a corresponding subjective one, since
$Y = 0$ gives $O = 0$, and that conversely every subjective moment in a
real objectivization must be determined by a corresponding objective one,
since $X = 0$ gives $O = 0$. Subjective and objective are, in the
ontological sense, correlative concepts, but not coordinate: the
objectivizing subjectivity is, as the overreaching unity of the opposition,
a \textit{dignitate prius}.

From the fundamental law it follows further that the truth of cognition
rests neither on the objective by itself nor on the subjective by itself,
but on the reciprocal determination of both. Let there be thousands,
indeed millions, of diversely constituted sense-systems in the universe,
still the a priori forms of cognition are the subjective fundamental forms
of universal reason, and universal reason is common to all cognizing
beings. The cognizing beings that, with the same sense-system, come under
the same
% p. 27
order must see that the concrete objectivization conditioned by sensation
has all the truth it can and should have, when the reciprocal action
between the subjective and the objective is precisely determined for the
order in question. The dissimilar sense-systems are just as little able to
awaken doubt about the universal validity in the peculiarity of
objectivization as the red-blind man, despite the truth of his
objectivization, is able to awaken doubt about the difference of colors
among those who have color-sense. The truth is that the reciprocal action
between subjective and objective must be different for subjectivities of
different order; but \emph{the reciprocal action itself is the truth}. To
assume an X standing outside the reciprocal action, unobjectivized, for the
really true, is a dream, an imagining; outside the reciprocal action there
is no truth, for outside the reciprocal action there is, ontologically
understood, nothing at all.


% ============================================================
%  II. THE PROBLEM OF REALITY
% ============================================================
\section*{II.\ The Problem of Reality}
\label{sec:realitet}

% p. 27 (heading printed here; section II text begins after it)
\noindent
While Hegelianism has sought to transform the world of realities into a
system of metaphysical concepts, Positivism has, for a change, thrown all
metaphysical ballast overboard, in order to hold exclusively to realities
and facts. It is, as is well known, not the realities that conform to our
concepts, but conversely our concepts that must conform to the realities;
the wisest thing philosophy can do is therefore to give up the a priori
problems, abandon its old leaky vessel, go on board the steamer of physics
as a passenger, and assume the title of Positivism. In justification of so
important a step, the school's celebrated founder, Auguste
Comte\footnote{\textit{Cours de Philosophie Positive; Tome 1mier, Paris
1830. 1miere leçon.}}, has pointed out three stages of development
(\textit{trois états}) which the human mind must necessarily pass through on
its way to the cognition of truth, \emph{the theological, the metaphysical,
and the scientific}. At the theological stage one seeks the cause of the
phenomena in a supernatural will and
% p. 28
lets the arbitrariness of the supernatural will determine the course of
things without regard to fixed laws. At the metaphysical stage the
arbitrariness disappears and gives place to such supersensible, self-valid
entities as the Absolute, Substance, the Cause, the Idea, and so on. This
intermediate stage has been necessary as a transitional step from the
theological to the scientific\footnote{\textit{La théologie et la physique
sont si profondément incompatibles, leurs conceptions ont un caractère si
radicalement opposé, qu'avant de renoncer aux unes pour employer
exclusivement les autres, l'intelligence humaine a dû se servir de
conceptions intermédiaires, d'un caractère bâtard, propres, par cela même, à
opérer graduellement la transition. Anfrt.\ Skr.\ S.\ 13.}}; but in itself
the metaphysical standpoint is barren, as can be seen particularly from the
endless, resultless disputes into which it dissolves. The positive
philosophy, which stands at the scientific stage, now characterizes itself
by this, that it regards the phenomena as subject to unchangeable laws,
seeks to reduce these laws to the smallest possible number of fundamental
laws, and considers every attempt to penetrate either to the first causes or
to the so-called final causes as utopian and
meaningless\footnote{\textit{Nous voyons, que le caractère fondamental de la
philosophie positive est de regarder tous les phénomènes comme assujétis à
des \emph{lois} naturelles invariables, dont la découverte préçise et la
réduction au moindre nombre possible sont le but de tous nos efforts, en
considérant comme absolument inaccessible et vide de sens pour nous la
recherche de ce qu'on appelle les \emph{causes}, soit premières, soit
finales. Anfr.\ Skr.\ S.\ 14. Further, S.\ 16: La preuve manifeste de
l'impossibilté d'obtenir de telles solutions, c'est que, toutes les fois
qu'on a cherché à dire à ce sujet quelque chose de vraiment rationnel, les
plus grands esprits n'ont pu que définir ces deux principes --- l'attraction
et la pesanteur --- l'un par l'autre, en disant, pour l'attraction, qu'elle
n'est autre chose qu'une pesanteur universelle, et ensuite, pour la
pesanteur, qu'elle consiste simplement dans l'attraction terrestre.}}.

To exclude metaphysics is fully justified when by the word metaphysics one
denotes a science that puts imagination in the place of experience in order
to discover the essences and \textit{qualitates occultæ} of which real
things must absolutely consist, unconcerned whether it agrees with reality
or not. But against so superficial a conception of the concept of
metaphysics one must protest in the name of science. Hegel's
``Wissenschaft der Logik'' is essentially metaphysics. Whatever imperfec-
% p. 29
tions may be pointed out in this Logic, one should nevertheless not accuse
the author of fetching his concepts from the world of imagination. The
concepts are so inherent in the things that no existing thing can escape
corresponding to its concept. That ``cultivation and discipline'' of
consciousness\footnote{``Das System der Logik ist das Reich der Schatten,
die Welt der einfachen Wesenheiten, von aller sinnlichen Konkretion befreit.
Das Studium dieser Wissenschaft, der Aufenthalt und die Arbeit in diesem
Schattenreich ist die absolute Bildung und Zucht des Bewusstseyns.'' Hegel.
Wissenschaft der Logik. Berlin 1833. Einleitung. S.\ 47.} which Hegel
demands on behalf of Logic he can with full right demand on behalf of all
the sciences. --- The relation between physics and metaphysics we may
determine thus: physics forms its concepts according to the observed facts,
metaphysics develops its concepts out of the conditions necessary for all
being and existence. To deduce the concretions of the concepts of existence
from pure thought is impossible; but to give an account of empirical
phenomena without going back to the a priori concepts is also impossible.
Gold, silver, iron are called partly different \emph{things}, partly
different \emph{substances}; what use is it now to denote reality, the given
fact, now by the word \emph{thing}, now again by the word \emph{substance},
when one is not in agreement with oneself as to wherein the designation
substance differs from the designation thing? In science one must after all
know what one is oneself saying. In order to be a \emph{definite thing}, the
phenomenon must absolutely be a \emph{thing}; in order to be a \emph{definite}
substance, it must be a \emph{substance}. On this point physics has the
advantage of consulting metaphysics. In metaphysics, the objective logic, it
is not only made evident that the thing must necessarily reflect itself in
the properties, and the properties in the thing, the substance in the
accidents and the accidents in the substance; but it is also made evident
wherein the reflection consists, and how different it is in the case of the
thing and in that of the substance.

From this standpoint it will be instructive to see two such strong thinkers
as Comte and Hegel try their strength with
% p. 30
one another. Comte maintains that there is no Absolute; and he is right
insofar as he looks only at existence and the existing. In the world of
facts the one term stands against the other, the one atom exists beside the
other. Everything is parceling-out, everything is relation. Hegel maintains
that the only thing that in truth is, is the Absolute; and he is right
insofar as he looks at the original unity of being with universal thought.
The relative is not merely relative in relation to another relative; but all
relations and relative terms have their subsistence in an encompassing
whole. If the encompassing whole is itself relative, it must necessarily
have its subsistence in relation to the relative term. But in this
thoroughgoing relation it then relates, through all relativities, to itself,
and precisely thereby comes to subsist as the Absolute. Comte maintains that
all that is logical is subjective; and he is right insofar as he thinks only
of formal logic, for formal logic with its abstractions is undeniably
subjective. Hegel maintains that the logical is essentially objective; and
he is right insofar as he starts from the essential unity of thinking and
being. The determination of thought to which no being corresponds is just
as empty as the determination of being to which no thought corresponds.

Comte assumes that thing and reality can very well be determined in
themselves and present themselves to us without any objectively determining
concept; but here he is wrong. We abstract, as it is said, our concepts from
the things; but to abstract concepts from things in which there are no
concepts is just as impossible as to pluck hairs from the bald man. Hegel
assumes that the objective concepts in their immediacy are themselves
things, and that it must therefore be possible to deduce existing facts from
logical universal concepts; but here he is wrong. How unhappily every
attempt at a construing of real things \textit{a priori} must turn out, his
philosophy of nature has more than sufficiently demonstrated.

Comte teaches that reality is reality, but not concept; Hegel teaches that
the concept as concept is reality: in the clearly
% p. 31
expressed opposition between Positivism and absolute Idealism the
\emph{problem of reality} can properly first be said to arise.

The problem is: that thing and concept must be assumed both to be the same
and yet not the same.

That thing and concept must, in an essential sense, be the same can easily
be seen. Take whatever thing or existing object you will, and you will
convince yourself that the sum of all determinations—be they material or
ideal, sensible or unsensible—must, since the determinations characterize
the object and make the thing what it is, necessarily belong to the object's
concept. In its sum of determinations the object subsists; with a change in
the determinations the object changes; with the dissolution of the
determinations the object disappears. When it nevertheless seems as though
there lay something more in the object itself than in its concept, this
derives simply from a confusion of the object's concept with our concept of
the object; for our concepts of concrete objects are far from being
exhaustive. That gold's color, malleability, specific gravity, form of
crystallization, slight affinity for oxygen, and so on, are properties that
characterize this metal and to that extent belong to the concept of gold,
there is surely no one who will deny. When we must nevertheless admit that
we have only an imperfect concept of gold itself, the reason for this is
easy to explain. We are, after all, still far from knowing all of gold's
properties, and even if we knew them from experience, we could still not
boast of knowing gold's substance so long as we are not able to derive the
properties from one another mutually or from the common unity in which they,
as accidents, have their subsistence. To be yellow, to be malleable, to
crystallize in octahedra, to have affinity for chlorine, cyanogen, and so
on, are in reality such dissimilar determinations that no mortal is able
either to point out their law-determined dependence on one another mutually
or to discover the real unity from which they all with necessity proceed.
And yet physics must concede that color, malleability, crystal-
% p. 32
form are just as much functions of the gold-atoms' quality and arrangement
as gold's unequal affinity for chlorine and oxygen; in other words: physics'
conception of the laws is not limited to certain relations between thing and
thing in outward generality; but every thing with its properties, every
whole with its parts, every substance with its accidents, every body with
its molecular structure expresses a system of law-determinations in which
term and relation reflect one another. Even the smallest parts can, in the
all-sided reciprocal action, just as little fall outside the reflection
comprehending its moments as they can fall outside the laws. Since now all
determinations, both of the particular terms and of the terms' mutual
relations, enter as real moments into the determining reflection, and the
totality of the reflected determinations is one with the concept, it follows
from this that thing and concept must be the same.

And yet concept and thing are by no means the same; they are, on the
contrary, completely opposed to one another. What a difference between the
concept of gold and the existing gold! The existing is immediately a fact,
an object of experience; but the concept as concept tolerates no object of
experience, no residues of immediacy, whether of sense-impressions or of
images or of representations: everything must be dissolved, melted, purified
in that fire of negativity which gives the concept its ideality's
unsensibleness and logical distinctiveness. From this standpoint Comte
undeniably has ground to make a divorce between concept and fact; for what
holds of the fact does not hold of the concept. One can crush a stone, but
not the stone's concept; one can boil a plant, but not the plant's concept;
one can shoot a bird, but not the bird's concept; the concept can be
thought, and one can think \emph{of} the thing, \emph{about} the thing, but
the thing as thing cannot be thought. The concept, the universal, can be
expressed in words, but the thing's individual determinateness must be
pointed at as at a This and This.

The immediacy that distinguishes thing from concept cannot be logical; for
even though the logical can point beyond itself, this always happens within
the sphere of logic: Hegel's
% p. 33
Logic is an eloquent example of this. Nor can the immediacy be illogical; if
the object-determinations were illogical and stood in contradiction with the
logical, the object would either be concept-less or stand in contradiction
with its concept. Nor can the immediate in the object be something that is
added or joined from outside to the logical, as when in certain logics a
distinction is made between the object itself and the concept expressed in
the object. What indeed should this something be? The material, matter? But
if matter with all its determinations falls outside the concept, how then
can the concept be expressed in the material object?

When now the sought immediacy can be neither logical nor illogical nor
joined to the logical in an outward manner, then it must surely be the
not-logical, which the logical needs, a necessary complement with which it
is saturated, in that it disappears; this complement is the antilogical.

If we seek the logical root from which the antilogical springs, we find it
in the category of power. Power in its immediacy is the expression of
reality, and the expression of reality's immediacy is again the antilogical.
The reason we regard the objects as realities and ascribe to them an
immediate, objective self-subsistence is precisely the impression of the
power with which they, each singly, come to meet us. The objects of
experience appear in this light like great and small potentates, who
maintain their place in proportion to the power they have; only that, for
the sake of immediacy, we must add that, without all logical reflecting,
they are what they have, and have what they are. In apprehending facts,
things, and realities from this standpoint, we have, in consequence of the
logical, come out beyond the logical. The opposition of thing and concept is
so striking that, when the thing is posited, the concept has disappeared,
and when the concept is posited, the thing has disappeared.

If we now go back to the fundamental law of ontological objectivization,
which was treated under the problem of knowledge, then it becomes possible
for us to point out the ground both of Positivism's
% p. 34
and of Hegelianism's one-sidedness. In that the Positivists immediately make
straight for facts and real objects, they entirely forget to examine the
relation between object and objectivization. The law says: as the object is,
its objectivization must be; as the objectivization is, the objectivizing
subjectivity must be. Since the adherents of the positive philosophy quite
overlook this law, they cannot, with all their interest in the exact
investigation of reality, of course get their eyes open either for
objectivization in power or for the objective real-concepts by which the
realities are \textit{a priori} determined. They are zealous for the
cognition of the unchangeable laws, without a suspicion that the
unchangeableness of the laws has its ground precisely in the system of
real-concepts. The best proof of this is the fact that Stuart Mill, whose
inductive logic has set empirical induction in the clearest light, himself
admits that the laws do not demonstrably belong to the essence of things,
that one may very well imagine regions of the heavens where ``events follow
one another by pure chance without any fixed law.'' Just as little have the
Hegelians, despite their absorption in the dialectic of the objective
concepts, been attentive to the fundamental law of objectivization. They
have categories such as ``Becoming,'' ``Appearing,'' ``Mediation,''
``Passing-over into its opposite,'' and the like; but the category of
objectivization is not to be found in the whole system. Had Hegel but once
fixed his sharp, penetrating gaze on the problem of objectivization, then we
should surely have got another explanation of the transition from the
subjective to the objective than is the case in that part of his Logic which
he calls the subjective\footnote{``Der Schluss ist \emph{Vermittelung}, der
vollständige Begriff in seinem \emph{Gesetztseyn}. Seine Bewegung ist das
Aufheben dieser Vermittelung, in welcher nichts an und für sich, sondern
jedes nur vermittelst eines Andern ist. Das Resultat ist daher eine
\emph{Unmittelbarkeit}, die durch \emph{Aufheben der Vermittelung}
hervorgegangen, ein \emph{Seyn}, das eben so sehr identisch mit der
Vermittelung und der Begriff ist, der aus und in seinem Andersseyn sich
selbst hergestellt hat. Dieses Seyn ist daher eine \emph{Sache, die an und
für sich ist}, --- die Objektivität.'' Hegel. Die subjective Logik. Berlin
1834. S.\ 170--171. When the subjective movement through mediation has
reached objectivity, and the immediate totality of being-determinations
stands before the logical consciousness as a factually given, it shows
itself what kind of subjectivity it was that comprehendingly found the
object. It is by no means the substance comprehended in the causal relation
that has transformed itself into judging, inferring, mediating subjectivity
and thereby produced the objective; it is quite simply a human reflecting
subjectivity, which has assured itself that there must be concept in the
world of objects, and now --- without in the least seeing how this concept
has become object, contenting itself with the assurance that the object is a
concept --- regards it from the outside, and seeks therein to find and
recognize the formal determinations that the subjective logic has developed.
Here there can be no question of ontological objectivization. --- ``Darum
ist das Erklären der Bestimmung eines Objekts, und das zu diesem Behufe
gemachte Fortgehen dieser Vorstellung nur ein \emph{leeres Wort}, weil in
dem andern Objekt, zu dem sie fortgeht, keine Selbstbestimmung liegt. Indem
nun die Bestimmtheit \emph{eines Objekts in einem andern liegt}, so ist
keine bestimmte Verschiedenheit zwischen ihnen vorhanden; die Bestimmtheit
ist nur \emph{doppelt}, einmal an dem einen, dann an dem andern Objekt, ein
schlechthin nur \emph{Identisches}, und die Erklärung oder das Begreifen
insofern \emph{tautologisch}. Diese Tautologie ist das äusserliche, leere
Hin- und Hergehen; da die Bestimmtheit von den dagegen gleichgültigen
Objekten keine eigenthümliche Unterschiedenheit erhält, und deswegen nur
identisch ist, ist nur \emph{Eine} Bestimmtheit vorhanden; und dass sie
doppelt sey, drückt eben diese Aeusserlichkeit und Nichtigkeit eines
Unterschiedes aus. Aber zugleich sind die Objekte \emph{selbstständig} gegen
einander; sie bleiben sich darum in jener Identität schlechthin
\emph{äusserlich}.'' Die subjektive Logik. S.\ 183--184. Consistently, the
development of the concept under this presupposition ought to have gone the
way of experience, to have penetrated the method of natural science and
thereby secured for itself the concepts' expression in facts; but instead
the illusion of the self-comprehending concept giving itself objectivity
gained the upper hand to such a degree that one, after having intoxicated
oneself in the dialectic of the immanent method, has now quite given up the
dialectical --- although dialectic is precisely a thoroughgoing critique of
ideality --- and is close to forgetting ``die absolute Bildung und Zucht des
Bewusstseyns,'' which philosophy can never dispense with, and which for so
essential a part is nevertheless owing to the great master.}. If it is
already from a logical standpoint evident that the object's concept cannot
possibly objectivize itself, that self-objectivization in the objectivization
of another is and must be subjectivity's ineffaceable characteristic mark,
then it follows irrefutably from this that a subjectivity which in thinking,
reflecting, and comprehending holds itself strictly within the boundary of
the logical will never be able to carry it further than to a logical
objectivization of logical objects.

% p. 36
In the real world the objects are not logical; they are real, immediate,
antilogical: their self-subsistence consists in immediate power, their
mutual relations are power-relations. What now lies in this? As the objects
are, the objectivization must be: power-objects must be objectivized in
power; further: as the objectivization is, the objectivizing subjectivity
must be: only an absolutely powerful subjectivity is able to let realities
proceed from their possibilities and bear a world of objects that really
stands by power. In the objects power is immediate as power in another; in
the objectivizing subjectivity power is reflected as self-power. Self-power,
the subjective power, is ideal and logical; the object-power, the objective
power, power in another, is real and antilogical. The opposition of ideal
and real power is the strongest, deepest opposition that can be thought.
Here is a real dualism; but the dualism is grounded in original unity, for
the absolute self-power has the opposition within itself and masters its
opposition.

The age-old question of the origin of the world-matter thereby comes under a
new light. Aristotle lets, as is well known, matter and form be equally
eternal, both belonging to that cycle of movements in which possibility
alternates with actuality, and actuality with possibility; matter is the
possibility, form the actuality. But this explanation is not appreciably
better than the more recent Materialists' assurance that matter and force
are equally eternal. What indeed is it, to be eternal, if it is not to be
the absolutely original? Absolutely original is, after all, only that which
has the ground of its origin in itself. But how is a being to be able to
have the ground of its origin in itself, if it does not at the same time
have a ground of origin for something other than itself? Aristotle describes
the first mover, which is not moved, the Godhead, as the being that does
nothing other than eternally to think the thinking of thinking (αὑτὸν ἄρα
νοεῖ \ldots{} καὶ ἔστιν ἡ νόησις νοήσεως νόησις)\footnote{Translator's note:
the Greek (Aristotle, \emph{Metaphysics} XII.9) reads roughly ``it thinks,
then, itself\ldots{} and its thinking is a thinking of thinking.''} and
which, in order not to stain its perfection with anything imperfect, must
exclude the possibility of being-other. But to exclude the possibility of
% p. 37
being-other is to exclude matter. That such a conception of the original
cannot withstand a critical test is easy to point out. If the eternal being
is to think itself, then this must necessarily happen by a movement of
thought in which it holds itself fast by distinguishing itself from something
other than itself; merely in order to think itself it needs an other. If,
moreover, its thinking is to be one with an active, real being, then this
active self-being must necessarily contain the possibility of being-other;
but the possibility of being-other is precisely matter. An eternal
self-being without matter-being, an eternal matter-being without self-being,
are hollow abstractions. The opposition of intellectual and real shows
itself here in its first ontological shade. One may now, with Spinoza,
attempt to lay the all-reality into the being's universal unity, conceive
the self-being as the infinite substance (\textit{causa sui}) with opposite
attributes, thought for the ideal and extension for the material, and reduce
the finite things to modifications, limitations, negations of the infinite:
it does not help. The abstraction sucks the blood of the substantial, so
that there remains neither self-being nor other-being, neither ideality nor
reality, that has any self-subsistence to show. The positivistic way out, to
draw a line through metaphysics and hold to the realities of experience, is,
as shown above, impassable; unless the decisive step is to be taken in all
thoughtlessness. For as regards the way out of thoughtlessness, there holds
of metaphysics what Hegel somewhere says of reason: the best way in which one
gets rid of it is: not to get involved with it at all.

If one is to hold fast reality without giving up thought, then the
ontological fundamental law holds, that the subjective, self-reflected ideal
must have the original power in order that the objective, the immediately
real, may have the derived. An equalization of the mighty opposition between
self-power and power in another cannot take place by subtraction and
addition. One does not increase self-power by drawing something from power
in another; nor can one strengthen the real by weakening or diminishing the
ideal. On the contrary! The greater the reality the
% p. 38
objective has, the mightier must the ideality be with which it is posited
and determined. The objective world of facts, of realities, is, in the
scientific sense, the strongest expression of the ideal power with which the
ontologically universal self-being masters and comprehends itself in all and
all in itself. All scientific truth stands and falls with the fundamental
principle that the finite subjectivity, which objectivizes in impotence,
must everywhere have its ontological equivalent in an infinite one, which
objectivizes in omnipotence.

The transition from ideal to real, from concept to thing, corresponding to
objectivization, has no intermediate determinations and can have none. How
indeed should these intermediate determinations be constituted? Logical they
cannot be, for one may enrich a concept with as many logical determinations
as one will, it still remains, as concept, within the logical. Antilogical
they cannot be either, for the antilogical enters only when the transition
has taken place. The transition must then be a power-movement, a conversion
of power's ideality into real power, a leap by which nothing is leapt over, a
\emph{logical-antilogical outburst of power}. Such an outburst of power is
originally one with the transition from \emph{possibility to
actuality}—understood, be it noted, in such a way that the possibility lies
in the concept, hence in the logical---\textit{potentia ab
posse}\footnote{Translator's note: the print here reads ``polentia af
posse'' (in roman), evidently a misprint for \textit{potentia ab posse},
``potency from being-able-to-be.''}---the actuality in the antilogical
reality, hence outside the logical. Let us see it in some examples. He who is
to walk with a light in a powder-magazine must be careful: he has an
explosion to fear. That the explosion is a transition from possibility to
actuality, an outburst of power, surely no one will deny. So long as the
powder remains in actuality, in antilogical reality—saltpetre, charcoal, and
sulfur in mechanical mixture—the explosion is only a possibility that lies
hidden in the powder's actuality; but when the explosion becomes actual, the
powder becomes unactual. The oxygen from the sulfate of potash unites with
the carbon into carbonic acid, the nitrogen develops in gaseous form, the
saltpetre's potassium forms with the sulfur potassium sulfide—in short: the
actual powder disappears; but the powder's concept, the real-concept powder,
remains behind in
% p. 39
the elements as possibility. The \emph{sudden} is here only a sensible
indication of the \emph{qualitative}. The possibility of hydrogen chloride,
a mixture of chlorine and hydrogen, can, by being placed in sunlight, show a
sudden turn into actuality. But the same transition can also take place
slowly, when the mixture stands in daylight; the transition can then be
regarded as a successive summation of manifold small, indeed infinitely
small, outbursts. The outburst rests on the qualitative: either the hydrogen
chloride's possibility or its actuality, in other words: either the ideal or
the real, either possibility or actuality, either concept or reality; an
intermediate there is not. To wish to condense the concepts into realities
or sublime the realities into concepts is a speculative bungling, in equal
degree contrary to experience and sound logic.

The whole real world rests upon an original outburst of power. If one
rejects the unity-power, the all-comprehending self-power, then the system
of real-concepts disappears and with it all the possibilities of reason. If
one assumes the self-power, then there follows with necessity, as lies in
its concept and comprehension in power, an outburst of antilogical
realities. The characteristic mark of the outburst's antilogical reality is,
seen from all sides, \emph{otherness}. The elements of actuality have a
double-sided stamp of otherness. That they are the other of self-power is
seen from this, that they constitute the real manifold, the swarm of atoms,
the parceling-out of immediate power through which the mighty unity confirms
itself. That there really is reality in this otherness is seen from this,
that there are objective relations, that the elements relate to one another
as other by other to other ad infinitum. We see it in the relation between
matter and force. Matter with its atoms is a striking expression of power in
another. It appears, on a superficial view, as though the atoms' existence
in space were only a collective coexistence with no other relation than that
which lies in the spatial relations. Consistently the Positivists must,
presumably, acquiesce in this assertion. By empirical induction one is
supposed to be able to convince oneself that the one atom-group follows
% p. 40
upon the other---usually!---according to a definite law. By taking motion as
a fact and appealing to the law, one gets free of force; that is, force is
transformed into a mathematical symbol, for example into the second
differential coefficient in accelerated motion. At this standpoint, however,
Materialism cannot maintain itself. The active relations between atoms and
atom-groups are too conspicuous for one to be able to escape the impression
of a real attraction and repulsion. The active relations now acquire a gleam
of reflection, and the reality of reflection makes itself known in force.
Then it is said that force, for example the force of gravity, is a property
of matter. But when force, the reflected-in-relation, is the active, then it
must surely also be the really actual, and what then is matter? One does not
rightly know whether it is matter that acts by virtue of its property,
force, or whether force is perhaps something particular, a kind of reflected
reality that adheres to matter, the immediate reality. The last assumption
is not satisfactory; so one goes further, draws the immediate reality under
reflection, makes matter itself into a phenomenon of forces, sublimes it
into forces and force-points. But---\textit{naturam furca expelles}---the
antilogical immediacy asserts its right; one is nevertheless compelled to
recognize a side of force that is \emph{opposite} to force, and to grant
matter, the immediate, self-less otherness, its inertia.

The unity-expression for matter and force is power in another, understood in
such a way that matter, the immediate in power, yields the existing terms,
and that force is reflected in the active relations. In matter the concept
founders, the existing stands in the concept's place; but that the
antilogical in existing matter is nevertheless logically determined by its
real-concept can be expressly known from force. That two masses existing in
space can attract one another at a distance is certainly incomprehensible
when one throws the concept away; the representation of an attractive force
adhering to matter gives no explanation. Only when one takes the
real-concept along does the
% p. 41
matter become intelligible. The existing masses are only one-sided
expressions of power in another. One supposes that one has reality's
complete expression in each of the masses by itself, and is astonished that
two such realities, which, as it seems, need nothing but the space they each
fill, can have anything to do with one another. That is an illusion. To
existing terms' actuality there belong active relations; the attractive
force, an expression of active relation, belongs inseparably in the
real-concept in which the masses' reality is originally grounded.

Let us pause a moment to hear an objection. From the standpoint of physics,
and then especially of Positivism, it must surely hold that attraction is a
fact, and that the discovery of the law of attraction, independent of any
analysis of metaphysical real-concepts, is owing solely to empirical
observation, calculation, and induction. To give the matter an appearance as
though that which depends exclusively on experience could also be
constructed \textit{a priori} is a barren shadow-boxing, a relapse into
antiquated metaphysics. This objection would be completely justified if it
did not have its ground in a misunderstanding. Here there is no question at
all of deriving real facts from our own imperfect, impotent concepts; on the
contrary! our task is to lead the facts of experience back to the
power-concepts, without which they, in their antilogical reality, would
stand before consciousness as blind facts. If it is in conflict with the
scientific method that one, with the help of experience, strives to derive
the laws of things from their concepts, how then will one defend the
procedure of mechanical physics? It is, after all, evident that the laws for
the impact of hard and elastic bodies cannot possibly be found by way of
experience alone, since no body is either perfectly hard or perfectly
elastic; they must necessarily be derived from the real-concepts: hardness
and elasticity.

From the ideal side the real-concepts are reflected in the all-comprehending
self-unity; from the real side they are inherent in the things. It is the
latter that in the present connection especially con-
% p. 42
cerns us. The actual contains possibility to another actual; the actual
gases oxygen and hydrogen contain possibility to water. In order that the
possibility may be actualized, conditions are required; the conditions have
the actual's immediacy and yet correspond in reflection to the possible.
Chemistry has examples in abundance. When the salt is decomposed, that is,
goes back into its concept, its possibility, the acid and the base step
forth in antilogical reality, and conversely; nothing but outbursts of power
in another. The real-concepts have a peculiar affinity for the antilogical;
this affinity is originally determined by the relation between power in
another and self-power. The chemist may concede that the oxygen and the
hydrogen in a certain way perish in the water, namely insofar as by the
combination they give up their particular properties, so that a new body
with new properties arises; but that this can be regarded as a return of the
realities to their real-concepts he will hardly concede. He will, in
consequence of the atomic theory, on the contrary reject the transformation
and assert the difference of the composite atoms from the simple atoms. The
water-atoms are decomposable, for each water-atom consists of one oxygen and
two hydrogen atoms; but the atoms of oxygen and hydrogen are, as simple
atoms, indecomposable. The ontological reversal from reality into concept,
from concept into reality, becomes by this mode of consideration highly
suspect. Yet wrongly. It is finitude's, otherness's, real discretion on
which the emphasis falls, but this real discretion is so far from conflicting
with the real-concept's logical essence that it on the contrary confirms it:
only by the discretion's antilogical elements entering as reflexes into
mighty comprehension does the concept acquire the character of real-concept.
The composite atom cannot possibly be regarded as a mere arithmetical sum of
two or more simple atoms; for the sake of aggregation alone a mutual
attraction must here be presupposed. Even if the attraction were—what, on
account of the peculiarity of the affinities, it cannot be—merely
mechanical, then the mechanical attraction, according to what we have seen
in the foregoing, would alone be sufficient to assert the concept's va-
% p. 43
lidity. The oxygen and the hydrogen could not possibly combine chemically if
they had no affinity for one another; but mutual affinity they could not
possibly have if their real-concepts did not correspond to one another; but
the real-concepts could not possibly correspond to one another if they were
not originally comprehended in one and the same comprehension. In the
oxygen, the hydrogen, and the water we have three real-concepts, two
co-concepts and a concretion-concept determined thereby. When the
co-concepts exist with antilogical reality, then the concretion-concept is a
possibility; when the concretion-concept exists, then we have in this
reality the co-concepts' possibility. The transition from possibility to
actuality presupposes conditions, but takes place by an outburst of power.

It has been intimated above that the positivistic assurance of the
immutability of the laws of nature stands, according to St. Mill's
admission, on rather weak feet. But even if this weakness is overlooked, even
if it is conceded that a merely empirical induction can give us a maximum of
probability that a B will invariably follow upon an A, is any essential
insight thereby gained? By no means. There is a difference between
\emph{certainty} and \emph{insight}. Suppose the relation between the
hypotenuse and the legs had been found practically by a measurement of the
triangle's sides and squares constructed thereon, then repeated measurements
of right-angled triangles could indeed give a certainty about equal to that
which a geometrical proof gives; but only the proof gives insight, the
measurement by itself gives no insight. All insight is conditioned by this,
that the representation of the one encloses the requirement of the other,
that the one cannot be posited without the other having to follow, that the
reality of the one bears within it the concept of the other. If the
hypothetical proposition ``if A occurs, then B must follow'' is held fast
merely empirically, then there is here no insight, however great the
certainty may be. One will ground the phenomenon by enumerating its
conditions, but however many conditions and intermediate determinations one
enumerates, one still does not reach the ground; for it stands with the
conditions as with the conditioned: the individual conditions are themselves
phenomena. Instead of a grounding of the
% p. 44
phenomenon one reaches only a summation of phenomena, which each in turn
must be grounded by a summation of phenomena; the grounding is a continued
demand for grounding, but the ground flees away into infinity: insight is
never reached. One can, with a maximum of probability, get certainty about
many things and about the connection of many things, without however
attaining insight, rational insight into any thing.

A rational insight into the connection of things is only possible under the
presupposition that there is reason in the things, that the things are
realized concepts. The concept realized in the thing is the thing's nature.
That every thing acts according to its nature means that every thing acts in
consequence of its concept; but what is the acting cause? For this question
the problem of concept and reality presents a new side. The concept seems
not to be able to be cause: the concept of gravity does not weigh, the
concept of fire does not burn, the concept of a stone kills no dog. It must
then be the thing that is cause. But when the thing can act only in
consequence of its concept, then the concept must surely be the cause of the
thing's being a cause. What does this lead to? Quite directly to a
distinction between two kinds of causes, ideal and real, and thereby to a
more penetrating investigation of the relation between \emph{ideal and real
causality} on the whole. Such an investigation is of importance for the sake
of insight and scientific cognition. To make matter and force into causes,
without troubling oneself further either about the concept of matter or about
the concept of force, is an easy matter; to make the real cause logically one
with the concept of the cause is not exactly difficult either; but in the
first case one excludes all insight, in the last all reality.

The logical-antilogical outburst of power, which transforms the concept into
reality, transforms at the same time the ideal causality into the real.

That the cause can act only by virtue of the unity of concept and reality is
seen from this, that acting and activity can take place only in consequence
of the same unity. The thing as thing in antilogical immediacy cannot act.
And why not? Because
% p. 45
immediacy, sharply and exactly held fast by itself, means that the concept's
content, and therewith possibility to another, is exhausted in the thing's
existence; the only thing the thing in this case is able to do is: to exist.
With exhausted concepts the things would exist side by side, excluding one
another, but for the rest have not the least to do with one another.
Herbart's absolutely positive reals can serve as an example of this. This
acute thinker, persistent in exact analyses, suffered from the fixed idea
that unconditionally independent reals should be able to be affected by one
another without the positive in them betraying anything negative. A curious
contradiction\footnote{The propensity to give empty verbal explanations
instead of a real interpretation of facts, and the weaknesses of which
metaphysicians, for the sake of the system alone, often make themselves
guilty, F. A. Lange (Geschichte des Materialismus und Kritik seiner
Bedeutung in der Gegenwart. Iserlohn. 1866. S.\ 463.) has aptly
characterized: ``Das absolut Einfache ist auch keiner inneren Veränderung
fähig, weil wir uns diese nur in der Form wechselnder Ordnung der Theile
denken können. Deshalb sagt Herbart auch nicht, dass die Realen auf einander
wirken, sondern dass sie Einwirkungen von einander erleiden \emph{würden},
wenn sie diesen nicht durch einen Akt der \emph{Selbsterhaltung} Widerstand
leisteten. Als ob damit etwas andres gesagt werden könnte, als mit der
Annahme einer einfachen Wechselwirkung! \emph{Waitz} legt in seiner
Psychologie (S.\ 81) viel Werth auf den Unterschied zwischen
\emph{Dispositionen} zu einem Zustand und wirklichen Zuständen. So geht es in
der Metaphysik. Zustände darf die Seele nicht haben; bei Leibe nicht, sonst
ginge ja ihre absolute Einheit verloren! Aber Dispositionen, das ist ganz
etwas andres; „Strebungen", warum nicht? Der Metaphysiker widerlegt mit
einem enormen Aufwand von Scharfsinn alle möglichen andern Ansichten, und wo
er seine eigne Meinung entwickelt, schiesst er einen logischen Purzelbaum von
der gewöhnlichsten Sorte. Jeder Andre sieht, dass eine Disposition zu einem
Zustand auch ein Zustand ist, dass Selbsterhaltung gegen eine
\emph{drohende} Einwirkung nicht ohne eine, wenn auch noch so feine
\emph{wirkliche} Einwirkung denkbar ist. Der Metaphysiker sieht dies nicht.
Er hat sich mit seiner Dialektik an den Rand des Abgrundes getrieben, alle
Begriffe hundertmal herumgewendet, hervorgezogen, weggeworfen, und endlich
\emph{muss durchaus und durchaus etwas gewusst werden}. Also die Augen
zugedrückt und den \emph{salto mortale} herzhaft gemacht --- von den Höhen
der schärfsten Kritik hinab in die allergewöhnlichste Verwechslung von Wort
und Begriff! Ist dies gelungen, dann geht es munter weiter. Je mehr
Widersprechendes in die erste Grundlegung aufgenommen wird, desto freier
lässt sich schliessen, wie man denn bekanntlich aus mathematischen Sätzen,
welche den Faktor Null in versteckter Weise enthalten, oft die merkwürdigsten
Dingen ableiten kann.'' It is of course self-evident that a metaphysician's
error can as little be charged to metaphysics as a mathematician's
miscalculation to mathematics. If the mathematician subtly defends his
errors, when they are clearly pointed out, instead of openly admitting them,
he does not thereby compromise mathematics, he compromises only himself. The
same holds of the metaphysician and metaphysics.}! Under the mutual
% p. 46
affecting the reals must precisely lose their absolutely positive character
and sink down to becoming relativities. That it can no more be the concepts
which, held fast by themselves in logical form, produce the things, is
self-evident. The concept of water can as little produce water as one can
drown in the concept of water. The thing's concept is its possibility, a
possibility that lies hidden in other things' actuality. Only by the actual
things having possibilities in themselves to other things are they able to
act. Of several freely existing masses each single one has in its actuality
the possibility of changing the others' place in space; it is this unity of
actuality and possibility that is expressed in the activity we call
mechanical attraction. Two pieces of wood cannot be rubbed against one
another without the possibility of the one having to correspond to the
actuality of the other; light-activity cannot be given unless the
light-giver's actuality corresponds to the ether's possibility, and
conversely. Since now actuality is reality, and possibility concept, it is
seen from this that ideal and real activity, ideal and real cause, must
presuppose one another mutually.

The unity of ideal and real causality is a unity in opposition, an
opposition-unity. From the side of reason the natural order of things makes
the impression of an infinite, in all parts coherent whole; not even the
smallest part can withdraw itself from the power of necessity, for
everything subsists in infinite reciprocal action; it is a world of inner,
essential relations, a world of concepts of ideal causality. From the
standpoint of the sensible experiences, on the other hand, the world of
realities shows itself as a world of chance, of outward conditions and real
causality. According to the connection of concepts there has, for example,
from a
% p. 47
combination between silicon (silicium) and oxygen arisen silicic acid, from
a combination between potassium and oxygen potash, from a combination
between aluminium and oxygen alumina, and from the mutual combination of
these combinations silicate of potash-alumina (feldspar)—all essentially in
consequence of ideal causality. The one-sided concepts are contained in the
more all-sided ones, and every trace of chance seems to have disappeared. By
analysis of the concretion-concept feldspar we come to the co-concepts
silicic acid, potash, and alumina; by analysis of these still concrete
co-concepts to the elementary concepts of oxygen, silicon, potassium,
aluminium—with reflected necessity on account of essential relations. If we
begin with the elementary concepts and apply synthesis, we get the same
necessity within the logical. But in the antilogical sphere of the realities
something quite different shows itself. Here we meet circumstances, chances,
and contingencies on account of outward relations. The oxygen and the
silicon can here exist each by itself, without silicic acid arising; silicic
acid and potash are found each by itself, without silicate of potash
forming, and so forth. In the circle of the outer conditions every outburst
of power is dependent on intervening chance. The outburst, antilogical in
the contingency, is in the necessity logical: the reciprocal determination
of contingent and necessary characterizes the law.

This leads us a step further. There is in the system of real-concepts a
difference between concept-values, marked by this, that some concepts
apparently lie nearest to necessity, others nearest to contingency. Concepts
such as earth, water, salt, acid, base, and the like appear as
essence-concepts, substance-concepts, concepts of necessity; while others,
such as: gravel, heap, aggregate, conglomerate, and so on, are formed in an
entirely contingent manner by the coincidence of concept-consequences that
do not essentially concern one another; such concepts are owing to
contingency to such a degree that they do not even seem to deserve the name
of concepts. This distinction between concepts of substantiality and of
accidentality has a claim to atten-
% p. 48
tion, for the solution of the problem of reality is dependent on how it is
interpreted.

If we begin from the contingent and concede that contingency makes the
concept nominal, then the concept gravel is only a denomination for a
heaping-together of parts, but no objective real-concept. The Positivist
will smile when he hears that the concept gravel too is to be included in the
eternal comprehension and have its ontological origin from an
all-comprehending self-power. But when the concept gravel is nominal, what
then is the concept granite? Also nominal! When the concept granite is
nominal, what then is the concept feldspar? Naturally also nominal; the
formation of feldspar has, after all, been conditioned by outward relations.
That this mode of consideration is false shines forth from its consequence;
for the consequence is that all objective real-concepts must disappear and
give place to nominal designations of concept-less facts. The absurdity and
superficiality herein has been sufficiently pointed out above; we are then
compelled to begin from necessity. Just as in all real necessity there is an
antilogical moment of contingency, since the really necessary is always
\emph{conditioned}, so there is in all the contingent a logical moment of
necessity; if it were not so, there would open up in existence a wide field
where the contingencies could carry on their play without regard to the
laws.

There have been carried on, in the history of the sciences, particularly of
philosophy, extensive disputes about whether, besides the physical causes,
one could also assume final causes and ascribe to these objective validity.
By distinguishing between the outer and the inner purposiveness, Kant has
held that, for the organic formations, one ought to make—not indeed a
constitutive, but nevertheless a regulative—use of the principle of
purposiveness; but even the regulative use seems more and more to sink in
credit. The obscurity has its ground in the fact that the relation between
ideal and real causality has not been sufficiently investigated.

When one without further ado separates the things from their concepts and,
with denial of the ideal causality, imagines that one
% p. 49
has enough in the physically acting causes, the final causes cannot possibly
acquire any objective significance. But that this tearing-loose of the real
from the ideal is unscientific will gradually become evident through a more
thorough treatment of the problems of knowledge and of reality. Both the law
of objectivization and the reciprocal determination of the logical and the
antilogical in the outburst of power prove that the real causality has its
ground in the ideal. When one one-sidedly clings to particulars, contents
oneself with dissection, and blindly makes the whole into a result of the
parts, then it is quite true that all talk of purpose and end must become
nominal and superficial. Every finite attempt to fix something as end and
something else as means founders on the infinite reciprocal action of
conditions and conditioned terms; the alleged higher sinks down to the level
of the supposedly lower, and conditioned necessity becomes sole ruler in the
realm of things and laws. If reason is to come into its right, if a
scientific insight is to be possible, then every approach to the solution of
the great problems must at all points point toward a development of the
parts out of the whole, toward the grounding of the real totality in the
ideal, toward a world-design with an all-comprehending subjectivity as its
author.

In the real-concept \emph{design} lies the unity of ideal and real
causality, in the concept of design purpose, and in the concept of purpose
again the concept of means. The purpose, originally seen, is that
requirement of subjectivity that the plan in the whole, the definite design,
shall be actualized through a reciprocal development of the parts among
themselves. The design itself is a system of designs: no part and no whole
can give itself design, but every part, every whole acts and must act by
virtue of its peculiar design. There must, in the name of reason, be a whole
in order that there can be parts, there must be a higher in order that there
can be a lower, there must be purpose in order that there can be means. The
unscientific arises only when the finite subjectivity presumes to survey the
plan, determine the aims, and apportion sun and wind between objective
purposes and means.

% p. 50
The fundamental question is how the subjective in reality relates to the
objective. This problem itself comes into existence with the organic
formations. What is a living organism, indeed what is life? Life is a
conditioned self-activity, and the product of this activity is the organism;
that is the shortest answer we can give here. The self-activity conditioned
by matter exhibits a self-working-out, though restless, equilibrium of ideal
and real causality. At the lower, elementary stage of development, the ideal
causality, the ontologically subjective, is—what we have often enough
emphasized—present everywhere, but so hidden in realities that the things
act as matter-subjects by their properties upon one another: masses,
molecules, atoms are real causes, physical otherness-causes on account of
the hidden self-cause. With life's organizing activity the relation is
reversed: the otherness-causes sink down to conditions and means for the
self-cause coming to light, the self-actualizing ideal subjectivity, which
in reciprocal action with the manifold of matter-subjects makes itself the
ruling center and shows what the purpose is. In an enigmatic manner the
double-sided causality separates into two relatively independent
causalities, two heterogeneous principles, the ideal and the real; this must
then lead the inquiry further.


% ============================================================
%  III. THE PROBLEM OF THE HETEROGENEOUS PRINCIPLES
% ============================================================
\section*{III.\ The Problem of the Heterogeneous Principles}
\label{sec:ueensartede}
% [text to be added: pp. 50--77]

\end{document}
